\documentclass[11pt,a4paper]{article}

\newcommand{\runninghead}{Study 1 — Temperature compensation, current instrument era}
% ---------------------------------------------------------------------------
% reportstyle.tex — shared preamble for the study reports under studies/.
%
% Kept deliberately inside the package set shipped with TeX Live "basic":
% no tcolorbox, no siunitx, no enumitem, no titlesec. Everything below
% compiles with a stock pdflatex installation.
%
% Usage, from a report directory one level below a study folder:
%     % ---------------------------------------------------------------------------
% reportstyle.tex — shared preamble for the study reports under studies/.
%
% Kept deliberately inside the package set shipped with TeX Live "basic":
% no tcolorbox, no siunitx, no enumitem, no titlesec. Everything below
% compiles with a stock pdflatex installation.
%
% Usage, from a report directory one level below a study folder:
%     % ---------------------------------------------------------------------------
% reportstyle.tex — shared preamble for the study reports under studies/.
%
% Kept deliberately inside the package set shipped with TeX Live "basic":
% no tcolorbox, no siunitx, no enumitem, no titlesec. Everything below
% compiles with a stock pdflatex installation.
%
% Usage, from a report directory one level below a study folder:
%     \input{../../_shared/reportstyle.tex}
%     \graphicspath{{../outputs/}}
% ---------------------------------------------------------------------------

\usepackage[T1]{fontenc}
\usepackage[utf8]{inputenc}
\usepackage{textcomp}
\usepackage{lmodern}
\usepackage{microtype}
\usepackage[margin=2.4cm,headheight=15pt]{geometry}
\usepackage{graphicx}
\usepackage{booktabs}
\usepackage{tabularx}
\usepackage{array}
\usepackage{amsmath}
\usepackage{xcolor}
\usepackage[font=small,labelfont=bf,justification=justified]{caption}
\usepackage{float}
\usepackage{fancyhdr}
\usepackage[hidelinks]{hyperref}

% --- Palette ---------------------------------------------------------------
\definecolor{rulegrey}{HTML}{B9C0C7}
\definecolor{fillgrey}{HTML}{F2F4F6}
\definecolor{failred}{HTML}{A5202B}
\definecolor{passgreen}{HTML}{1B6B3A}
\definecolor{accent}{HTML}{1F4E79}

% --- Semantic shorthands ---------------------------------------------------
\newcommand{\degC}{\textdegree C}
\newcommand{\mdegC}{mdeg/\textdegree C}
\newcommand{\fail}{\textcolor{failred}{\textbf{fail}}}
\newcommand{\pass}{\textcolor{passgreen}{\textbf{pass}}}
\newcommand{\worse}{\textcolor{failred}{worse}}
\newcommand{\better}{\textcolor{passgreen}{better}}

% --- Highlighted panel -----------------------------------------------------
% #1 = heading, #2 = body
\newcommand{\panel}[2]{%
  \begin{center}%
  \setlength{\fboxsep}{9pt}%
  \fcolorbox{rulegrey}{fillgrey}{%
    \begin{minipage}{0.93\textwidth}%
      {\bfseries\color{accent}#1}\par\medskip
      #2
    \end{minipage}}%
  \end{center}}

% --- Numeric column: right aligned, fixed width ----------------------------
\newcolumntype{R}[1]{>{\raggedleft\arraybackslash}p{#1}}
\newcolumntype{L}[1]{>{\raggedright\arraybackslash}p{#1}}
\newcolumntype{Y}{>{\raggedright\arraybackslash}X}

% --- Table look ------------------------------------------------------------
\renewcommand{\arraystretch}{1.15}
\setlength{\tabcolsep}{6pt}

% --- Headers ---------------------------------------------------------------
\pagestyle{fancy}
\fancyhf{}
\fancyhead[L]{\small\itshape\runninghead}
\fancyhead[R]{\small\thepage}
\renewcommand{\headrulewidth}{0.4pt}

% --- Section spacing without titlesec --------------------------------------
\makeatletter
\renewcommand\section{\@startsection{section}{1}{\z@}%
  {-3.2ex \@plus -1ex \@minus -.2ex}{2.0ex \@plus.2ex}%
  {\normalfont\Large\bfseries\color{accent}}}
\renewcommand\subsection{\@startsection{subsection}{2}{\z@}%
  {-2.6ex \@plus -1ex \@minus -.2ex}{1.4ex \@plus.2ex}%
  {\normalfont\large\bfseries}}
\makeatother

% --- Figure defaults -------------------------------------------------------
\setkeys{Gin}{width=\linewidth}
\captionsetup[figure]{skip=6pt}

%     \graphicspath{{../outputs/}}
% ---------------------------------------------------------------------------

\usepackage[T1]{fontenc}
\usepackage[utf8]{inputenc}
\usepackage{textcomp}
\usepackage{lmodern}
\usepackage{microtype}
\usepackage[margin=2.4cm,headheight=15pt]{geometry}
\usepackage{graphicx}
\usepackage{booktabs}
\usepackage{tabularx}
\usepackage{array}
\usepackage{amsmath}
\usepackage{xcolor}
\usepackage[font=small,labelfont=bf,justification=justified]{caption}
\usepackage{float}
\usepackage{fancyhdr}
\usepackage[hidelinks]{hyperref}

% --- Palette ---------------------------------------------------------------
\definecolor{rulegrey}{HTML}{B9C0C7}
\definecolor{fillgrey}{HTML}{F2F4F6}
\definecolor{failred}{HTML}{A5202B}
\definecolor{passgreen}{HTML}{1B6B3A}
\definecolor{accent}{HTML}{1F4E79}

% --- Semantic shorthands ---------------------------------------------------
\newcommand{\degC}{\textdegree C}
\newcommand{\mdegC}{mdeg/\textdegree C}
\newcommand{\fail}{\textcolor{failred}{\textbf{fail}}}
\newcommand{\pass}{\textcolor{passgreen}{\textbf{pass}}}
\newcommand{\worse}{\textcolor{failred}{worse}}
\newcommand{\better}{\textcolor{passgreen}{better}}

% --- Highlighted panel -----------------------------------------------------
% #1 = heading, #2 = body
\newcommand{\panel}[2]{%
  \begin{center}%
  \setlength{\fboxsep}{9pt}%
  \fcolorbox{rulegrey}{fillgrey}{%
    \begin{minipage}{0.93\textwidth}%
      {\bfseries\color{accent}#1}\par\medskip
      #2
    \end{minipage}}%
  \end{center}}

% --- Numeric column: right aligned, fixed width ----------------------------
\newcolumntype{R}[1]{>{\raggedleft\arraybackslash}p{#1}}
\newcolumntype{L}[1]{>{\raggedright\arraybackslash}p{#1}}
\newcolumntype{Y}{>{\raggedright\arraybackslash}X}

% --- Table look ------------------------------------------------------------
\renewcommand{\arraystretch}{1.15}
\setlength{\tabcolsep}{6pt}

% --- Headers ---------------------------------------------------------------
\pagestyle{fancy}
\fancyhf{}
\fancyhead[L]{\small\itshape\runninghead}
\fancyhead[R]{\small\thepage}
\renewcommand{\headrulewidth}{0.4pt}

% --- Section spacing without titlesec --------------------------------------
\makeatletter
\renewcommand\section{\@startsection{section}{1}{\z@}%
  {-3.2ex \@plus -1ex \@minus -.2ex}{2.0ex \@plus.2ex}%
  {\normalfont\Large\bfseries\color{accent}}}
\renewcommand\subsection{\@startsection{subsection}{2}{\z@}%
  {-2.6ex \@plus -1ex \@minus -.2ex}{1.4ex \@plus.2ex}%
  {\normalfont\large\bfseries}}
\makeatother

% --- Figure defaults -------------------------------------------------------
\setkeys{Gin}{width=\linewidth}
\captionsetup[figure]{skip=6pt}

%     \graphicspath{{../outputs/}}
% ---------------------------------------------------------------------------

\usepackage[T1]{fontenc}
\usepackage[utf8]{inputenc}
\usepackage{textcomp}
\usepackage{lmodern}
\usepackage{microtype}
\usepackage[margin=2.4cm,headheight=15pt]{geometry}
\usepackage{graphicx}
\usepackage{booktabs}
\usepackage{tabularx}
\usepackage{array}
\usepackage{amsmath}
\usepackage{xcolor}
\usepackage[font=small,labelfont=bf,justification=justified]{caption}
\usepackage{float}
\usepackage{fancyhdr}
\usepackage[hidelinks]{hyperref}

% --- Palette ---------------------------------------------------------------
\definecolor{rulegrey}{HTML}{B9C0C7}
\definecolor{fillgrey}{HTML}{F2F4F6}
\definecolor{failred}{HTML}{A5202B}
\definecolor{passgreen}{HTML}{1B6B3A}
\definecolor{accent}{HTML}{1F4E79}

% --- Semantic shorthands ---------------------------------------------------
\newcommand{\degC}{\textdegree C}
\newcommand{\mdegC}{mdeg/\textdegree C}
\newcommand{\fail}{\textcolor{failred}{\textbf{fail}}}
\newcommand{\pass}{\textcolor{passgreen}{\textbf{pass}}}
\newcommand{\worse}{\textcolor{failred}{worse}}
\newcommand{\better}{\textcolor{passgreen}{better}}

% --- Highlighted panel -----------------------------------------------------
% #1 = heading, #2 = body
\newcommand{\panel}[2]{%
  \begin{center}%
  \setlength{\fboxsep}{9pt}%
  \fcolorbox{rulegrey}{fillgrey}{%
    \begin{minipage}{0.93\textwidth}%
      {\bfseries\color{accent}#1}\par\medskip
      #2
    \end{minipage}}%
  \end{center}}

% --- Numeric column: right aligned, fixed width ----------------------------
\newcolumntype{R}[1]{>{\raggedleft\arraybackslash}p{#1}}
\newcolumntype{L}[1]{>{\raggedright\arraybackslash}p{#1}}
\newcolumntype{Y}{>{\raggedright\arraybackslash}X}

% --- Table look ------------------------------------------------------------
\renewcommand{\arraystretch}{1.15}
\setlength{\tabcolsep}{6pt}

% --- Headers ---------------------------------------------------------------
\pagestyle{fancy}
\fancyhf{}
\fancyhead[L]{\small\itshape\runninghead}
\fancyhead[R]{\small\thepage}
\renewcommand{\headrulewidth}{0.4pt}

% --- Section spacing without titlesec --------------------------------------
\makeatletter
\renewcommand\section{\@startsection{section}{1}{\z@}%
  {-3.2ex \@plus -1ex \@minus -.2ex}{2.0ex \@plus.2ex}%
  {\normalfont\Large\bfseries\color{accent}}}
\renewcommand\subsection{\@startsection{subsection}{2}{\z@}%
  {-2.6ex \@plus -1ex \@minus -.2ex}{1.4ex \@plus.2ex}%
  {\normalfont\large\bfseries}}
\makeatother

% --- Figure defaults -------------------------------------------------------
\setkeys{Gin}{width=\linewidth}
\captionsetup[figure]{skip=6pt}

\graphicspath{{../outputs/}}

\title{\vspace{-1.4cm}\bfseries Is the documented temperature compensation good for
the prediction system?\\[4pt]
\large\mdseries Study 1 --- the current instrument era, station 02\\
\large\mdseries Medieval walls of Gubbio, static structural health monitoring}
\author{}
\date{Notebook executed 11 August 2026 \quad$\cdot$\quad
Data window 21 February 2025 -- 10 August 2026 \quad$\cdot$\quad
Hourly analysis grid}

\begin{document}
\maketitle
\thispagestyle{fancy}

%% ==========================================================================
\section*{Summary}
\addcontentsline{toc}{section}{Summary}

The preprocessing stage of the pipeline applies a fixed thermal correction to the
inclinometer signal before any other stage sees it:
%
\begin{equation}
I_{\text{comp}}(t) \;=\; I(t) \;-\; \bigl(T_{\text{air}}(t) - T_{\text{air}}(t_0)\bigr)
\cdot c \cdot 1000 ,
\qquad c = 0.005 ,
\label{eq:comp}
\end{equation}
%
that is, five millidegrees of inclination removed per degree Celsius of air
temperature, applied instantaneously, driven by the station's own air-temperature
channel. This report asks two questions. Does that correction help or hurt a prediction
system built on this signal? And is its central premise --- that the structure responds
to temperature at the instant it is measured --- true?

\panel{Verdict on the correction: it fails all three criteria}{%
\begin{tabularx}{\linewidth}{@{}Y R{4.3cm} l@{}}
\toprule
\textbf{Criterion} & \textbf{Measured} & \textbf{Result}\\
\midrule
Thermal dependence is reduced
  & $|r|$: $0.845 \rightarrow 0.956$
  & \fail\\
Coefficient is near the variance-minimising value
  & $5.00$ vs $1.60$ \mdegC
  & \fail\\
Forecast improves under a controlled contrast
  & penalty $+68.2\,\%$ to $+119.3\,\%$
  & \fail\\
\bottomrule
\end{tabularx}%
\medskip

\noindent The measured thermal slope of this instrument is $1.63$ \mdegC\ (95\,\%
confidence interval $1.49$--$1.77$). The documented coefficient is roughly three times
that value, which places it on the far side of the variance minimum: applying it
inflates the residual variance to $3.3\times$ that of the untouched signal and
$11.7\times$ the achievable minimum, and it leaves the corrected series \emph{more}
strongly related to temperature than the raw series was --- with the sign reversed.}

\panel{Verdict on the premise: instantaneous in air temperature, lagged in radiation}{%
The response to \textbf{solar radiation} is measurably delayed. Filtering the
radiometer through a first-order thermal lag lifts the variance it explains from
$0.575$ to $0.731$ at a time constant of two hours, and from $0.346$ to $0.460$ at four
hours on the unfiltered signal. That is the physical mechanism working as expected:
flux arrives, a thermal mass integrates it, the structure moves afterwards.\par\medskip
The response to \textbf{air temperature} shows no delay at any timescale this window
can resolve. The best time constant is zero, and at the diurnal scale instantaneous air
temperature explains $93.2\,\%$ of the variance of the inclination. This is not
evidence against the mechanism --- it is what the mechanism predicts, because air
temperature is itself an already-integrated response to the same radiation and so
arrives pre-lagged. Air temperature works as a proxy for the forcing precisely because
it is a lagged one.\par\medskip
\textbf{Neither finding rescues the correction}, which fails on magnitude, not on
timing. Window A covers half an annual cycle; the seasonal response is a separate
question, and the companion legacy study answers it very differently --- there the
inclination \emph{leads} air temperature by about two months, which no thermal lag can
reproduce.}

\noindent\textbf{Recommendation.} Remove the fixed compensation from Notebook 00 and
let the grey-box decomposition estimate the thermal response. If a compensated series
is required for presentation, the coefficient must be fitted per station, never fixed.

\medskip
{\small\tableofcontents}

\clearpage
%% ==========================================================================
\section{The question and how it is decided}

\subsection{What is under test}

Equation~\eqref{eq:comp} is the correction as implemented in
\texttt{heritageshm/preprocessing.py} and applied in Notebook 00. Three of its
properties matter for what follows.

\begin{description}
\item[The coefficient is fixed.] $c = 0.005$ is a constant in the source, not a fitted
quantity, and it is applied identically to every station.
\item[The response is assumed instantaneous.] The correction uses the air temperature
recorded at the same timestamp as the inclination. It contains no lag term. This is an
assumption about the structure, and Section~\ref{sec:lag} tests it rather than
accepting it.
\item[The reference point is arbitrary.] $T_{\text{air}}(t_0)$ is the first retained
record, and the result is normalised so that the first value is zero. Both are
constants and both vanish under first differencing, so neither affects any conclusion
below.
\end{description}

\subsection{The physical hypothesis}
\label{sec:hypothesis}

The wall stands in an embankment condition: one face is exposed, the other is
restrained by the fill behind it. Solar radiation heats the exposed face, the heat
diffuses into a large masonry mass, and the differential expansion between the two
faces tilts the structure. Two consequences follow, and both are testable.

First, \textbf{the true forcing is radiation, not air temperature}. Air temperature is
a proxy --- the only one available wherever the radiometer is not, which is the entire
legacy era and every station but this one.

Second, \textbf{a body with thermal mass cannot respond instantaneously}. Whatever the
correction assumes, the structure must integrate its forcing over some characteristic
time. That does not by itself mean a correction driven by air temperature must carry a
lag term, because a proxy can already contain the delay. It means the delay has to be
measured in whichever variable is actually used.

\subsection{The three criteria for the correction}

A correction that claims to remove a thermal component makes falsifiable promises. The
study evaluates three, each decided by a number computed in the notebook rather than by
judgement.

\begin{enumerate}
\item \textbf{Thermal dependence must fall.} After correction, the absolute correlation
between the signal and air temperature must be smaller than before. This is the
correction's own stated goal and the weakest test it could be asked to pass.
\item \textbf{The coefficient must be near the variance-minimising value.} Sweeping $c$
traces residual variance as a function of the correction applied. A defensible
coefficient sits near the minimum of that curve; one past the minimum makes the series
worse than leaving it alone.
\item \textbf{The forecast must improve under a controlled contrast.} Two models of
identical form, differing only in whether the target was compensated, must not be
harder to forecast in the compensated case. This is the criterion that matters for the
pipeline, because the pipeline exists to forecast.
\end{enumerate}

\subsection{Prior evidence from the manufacturer documentation}
\label{sec:manufacturer}

A separate review of \texttt{docs/docs-manufacturer/} established the context in which
these results should be read. The instrument is a SEIKA SB1U capacitive inclinometer
with an integrated conditioner: zero point $2.5$~V, sensitivity approximately
$1000$~mV per degree, range $\pm 2$\textdegree. The logged \texttt{I} channel is the
conditioner output in millivolts, so one millivolt corresponds to one millidegree once
the $2500$ offset is subtracted.

Two statements bear directly on Equation~\eqref{eq:comp}. First, the datasheet lists
temperature-drift compensation as a \emph{built-in} feature of the conditioner, so a
software correction risks being applied on top of a hardware one. Second, the residual
drift specification --- 20~ppm per kelvin on the reference voltage --- bounds the
instrument's own thermal error at the order of $0.1$ \mdegC. The value applied in
software is $5$ \mdegC, roughly fifty times larger.

This also disposes of an alternative reading of the results below. The signal follows
air temperature closely and without delay, which is the signature a residual instrument
sensitivity would leave. It cannot be one: the measured slope is $1.63$ \mdegC, some
sixteen times the manufacturer's bound. The thermal response being measured is the
structure's, not the sensor's --- which is exactly why it belongs in the model rather
than in a preprocessing subtraction.

%% ==========================================================================
\section{Data and method}

\subsection{Sampling: the hourly grid}

Every environmental proxy the pipeline aligns against --- ERA5-Land and the other
reanalysis products --- is published hourly. The record is therefore aggregated from
its native 20-minute rate to one hour immediately after parsing, and every number in
this report is computed on that grid. Analysing at 20 minutes and resampling afterwards
would make the conclusions depend on a resolution the production pipeline never sees.

Each hour is the mean of its raw samples, kept only if at least two of the three
survived cleaning, so an hour rebuilt from a single reading is not presented as
equivalent to a complete one. Aggregation preserves the statistics that matter and
moderates the extremes, as it should: mean air temperature is unchanged at
$16.01$\,\degC\ and mean inclination at $-318.8$ mdeg, while the largest one-off solar
reading falls from $1387.8$ to $1358.0$~W/m². Of 38\,593 native slots, 12\,865 hourly
intervals result, 6 of them rejected for holding too few raw samples.

One capability is lost and is restored explicitly: an hourly grid cannot resolve a lag
shorter than an hour, so the native record is retained and scanned once, in
Section~\ref{sec:lag}, as a sub-hourly check.

\subsection{Station, era and windows}

Station 02, current instrument era --- the only period in which wall temperature and
solar radiation exist alongside the inclinometer, and therefore the only period in
which the question can be asked with every candidate driver present.

\begin{table}[H]
\centering
\caption{Channel coverage over the loaded span (12\,865 hourly intervals).}
\begin{tabular}{@{}l R{1.9cm} R{2.0cm} R{1.9cm} R{1.8cm} R{1.8cm} R{1.8cm}@{}}
\toprule
\textbf{Channel} & \textbf{Present} & \textbf{Available} & \textbf{Min} & \textbf{Max}
& \textbf{Mean} & \textbf{Std}\\
\midrule
\texttt{batt}  & 9\,919 & 77.1\,\% & 3.271 & 4.964 & 4.161 & 0.396\\
\texttt{tair}  & 9\,919 & 77.1\,\% & $-5.17$ & 38.32 & 16.009 & 8.924\\
\texttt{rh}    & 9\,919 & 77.1\,\% & 25.25 & 100.00 & 75.621 & 20.539\\
\texttt{sr}    & 8\,884 & 69.1\,\% & 0.00 & 1358.00 & 244.410 & 335.320\\
\texttt{twall} & 4\,706 & 36.6\,\% & 1.25 & 55.56 & 23.154 & 11.617\\
\texttt{inc}   & 9\,919 & 77.1\,\% & $-1712.1$ & $-261.6$ & $-318.8$ & 48.751\\
\bottomrule
\end{tabular}
\end{table}

Wall temperature is the binding constraint: its probe returns the $-55.0$ failure
sentinel across 15\,566 native records, leaving it valid in only $36.6\,\%$ of the
span. Two contiguous windows carry every channel simultaneously.

\begin{table}[H]
\centering
\caption{The two windows in which the inclinometer, both temperatures and solar
radiation are valid together.}
\small
\begin{tabular}{@{}l l R{1.4cm} R{2.8cm} L{5.0cm}@{}}
\toprule
\textbf{Window} & \textbf{Period} & \textbf{Days} & \textbf{Complete hours} &
\textbf{Role}\\
\midrule
A & 2025-02-21 to 2025-07-13 & 143 & 3\,415 of 3\,432 (99.5\,\%) & descriptive analysis and main experiment\\
B & 2026-06-18 to 2026-08-10 &  54 & 1\,225 of 1\,296 (94.5\,\%) & held-back out-of-period check\\
\bottomrule
\end{tabular}
\end{table}

Window B sits eleven months after Window A, in a different season, on the far side of a
103-day outage. It is never used to choose anything; it exists only to show whether the
ordering of the configurations survives a change of period.

\subsection{Signal conditioning}
\label{sec:screen}

The loader implements the contract in \texttt{docs/raw-data-format.md}: per-field
decimal normalisation, dates taken from the record rather than the filename, field-wise
merging of duplicate timestamps, every sentinel mapped to \texttt{NaN}, and the
inclinometer converted from conditioner millivolts to millidegrees about its calibrated
zero of 2500.

A gross-outlier screen then removes the instrument-settling transients that follow a
long outage, using a rolling-median context test with the threshold set from the robust
spread of the signal itself: $\sigma = 14.70$ mdeg, threshold $117.6$ mdeg, 59 of
9\,860 hourly samples removed ($0.60\,\%$).

\subsection{The six configurations}

\begin{table}[H]
\centering
\caption{The configurations compared under a common model. C against D is the
controlled contrast; E and F test the physics rather than the correction.}
\small
\begin{tabular}{@{}c l l L{6.0cm}@{}}
\toprule
& \textbf{Target} & \textbf{Regressors} & \textbf{Question it answers}\\
\midrule
A & raw & instantaneous & does the model benefit from the drivers directly?\\
B & compensated & instantaneous & what does correcting \emph{and} supplying the drivers cost?\\
C & compensated & none & the controlled contrast, compensated arm\\
D & raw & none & the controlled contrast, raw arm\\
E & raw & lag-optimised & does shifting each driver to its best lag help?\\
F & raw & inertia-filtered & does a first-order thermal lag help?\\
\bottomrule
\end{tabular}
\end{table}

Configurations E and F are built from quantities measured in Section~\ref{sec:lag},
not chosen: E shifts each driver by the lag at which its levels cross-correlation
peaks, and F filters each driver through the time constant that won the band-limited
inertia scan. Both act on the raw target, so they compete with A and D rather than with
the correction.

\medskip
\noindent\textbf{Every configuration is scored on the raw signal scale.} Predictions of
a compensated target have the correction added back before the error is computed, so
all six predict the same physical quantity. Without that step the comparison would be
meaningless, since a smaller mean absolute error could be bought by predicting a
flatter series. The notebook asserts that the compensation is exactly invertible before
any score is computed.

\medskip
\noindent\textbf{The controlled contrast is C against D.} Window A runs from late
winter to midsummer, so its test period sits largely outside the training temperature
range and any configuration that leans on temperature must extrapolate. Comparisons
\emph{between} regressor sets are therefore partly a measure of that extrapolation. C
and D share the same model form and the same absence of regressors, differing only in
whether the target was compensated. The verdict rests on them.

\medskip
The reference model is ridge regression on 12 autoregressive lags (12 hours of
history) plus the regressors, standardised, penalty $\alpha = 1$, fitted as a direct
forecast at the stated horizon, with a chronological 70/30 split.
Section~\ref{sec:np} repeats the comparison under NeuralProphet, the model the project
deploys, with a fixed random seed.

%% ==========================================================================
\section{What the signal and its drivers look like}

\begin{figure}[H]
\includegraphics{S1_02_st02_overview_windowA.png}
\caption{Window A --- inclination and every candidate driver on a common time axis,
21 February to 13 July 2025.}
\label{fig:ovA}
\end{figure}

The inclination carries a strong diurnal oscillation on top of a slow seasonal rise
that follows the advance of spring into summer. Nothing in the record suggests a
structural event: the series is environmental almost everywhere.

\begin{figure}[H]
\includegraphics{S1_03_st02_overview_windowB.png}
\caption{Window B --- the held-back period, 18 June to 10 August 2026, on the far side
of a 103-day outage.}
\label{fig:ovB}
\end{figure}

\subsection{Diurnal structure}

\begin{figure}[H]
\includegraphics{S1_04_st02_diurnal.png}
\caption{Mean diurnal profiles over Window A. Left, native units; right, standardised,
which is the informative panel because it compares shape and phase rather than
amplitude.}
\label{fig:diurnal}
\end{figure}

The standardised panel reads the phase relationships directly. The inclination and
air-temperature curves rise together from about 07:00, peak together near midday and
decay together through the afternoon: to the resolution of this plot they are in phase.
Wall temperature behaves differently, peaking near 16:00 --- the masonry probe lags the
air, as it must. Solar radiation leads slightly, peaking around 12:00--13:00 with the
steep morning rise characteristic of a clear-sky mean.

That the tilt tracks the \emph{air} rather than the \emph{wall probe} is worth holding
on to: it says the element governing the tilt responds faster than the mass the probe
is buried in. Section~\ref{sec:lag} makes this quantitative.

\subsection{Correlation structure}

\begin{table}[H]
\centering
\caption{Pearson correlation of the inclination with each driver over Window A, on the
hourly grid.}
\begin{tabular}{@{}l R{3.2cm} R{3.6cm}@{}}
\toprule
\textbf{Driver} & \textbf{Levels} & \textbf{First differences}\\
\midrule
\texttt{tair}  & $+0.845$ & $+0.826$\\
\texttt{twall} & $+0.791$ & $+0.394$\\
\texttt{sr}    & $+0.588$ & $+0.398$\\
\texttt{rh}    & $-0.611$ & $-0.673$\\
\bottomrule
\end{tabular}
\end{table}

\begin{figure}[H]
\includegraphics{S1_07_st02_corr_heatmaps.png}
\caption{Correlation structure over Window A.}
\label{fig:heat}
\end{figure}

Air temperature is the strongest driver on both views, which justifies the correction's
choice of channel even as the rest of this report rejects its magnitude. The drivers
are also strongly correlated with each other --- \texttt{tair} with \texttt{twall} at
$0.924$ in levels --- so they cannot be treated as independent explanations.

The differenced correlation with \texttt{tair} is $0.826$ here, against $0.375$ when
the same quantity is computed at 20-minute sampling. That is not a contradiction but a
property of differencing: at 20 minutes the difference of a smooth signal is dominated
by measurement noise, while an hourly difference is a cleaner estimate of the same
physical rate of change. It is also a warning about how strongly a differenced
statistic depends on the interval chosen --- and Section~\ref{sec:lag} takes that
warning seriously.

%% ==========================================================================
\section{Is the response instantaneous?}
\label{sec:lag}

\panel{A correction to an earlier version of this report}{%
An earlier version concluded, from a single cross-correlation of \emph{first
differences at 20-minute sampling}, that ``the instantaneous form of the correction is
defensible''. That inference does not hold. Differencing at the sampling interval is a
high-pass filter: it removes the slow component of both series by construction, so a
lag-zero peak in the differenced cross-correlation shows only that the fast wiggles are
synchronous. It cannot see a response governed by thermal mass, which is exactly the
response in question. This section replaces that single test with four: cross-correlation
on levels as well as differences, a sub-hourly check, and a scan over thermal time
constants on both the full signal and the diurnal band.}

\subsection{Lag structure on levels and on differences}

\begin{figure}[H]
\includegraphics{S1_09_st02_cross_correlation.png}
\caption{Cross-correlation of the inclination against each driver over $\pm 120$ hours,
on levels (left) and on first differences (right).}
\label{fig:ccf}
\end{figure}

\begin{table}[H]
\centering
\caption{Peak of each cross-correlation function. The levels search is restricted to
lags at which the driver leads and whose correlation carries the sign of the
contemporaneous one; without that restriction the near-periodic cross-correlation of
two cyclic series returns its anti-phase extremum, which is a property of periodicity
and not a lag.}
\begin{tabular}{@{}l R{2.4cm} R{2.2cm} R{2.6cm} R{2.2cm}@{}}
\toprule
& \multicolumn{2}{c}{\textbf{Levels}} & \multicolumn{2}{c}{\textbf{First differences}}\\
\cmidrule(lr){2-3}\cmidrule(lr){4-5}
\textbf{Driver} & \textbf{Peak $r$} & \textbf{Lag} & \textbf{Peak $r$} & \textbf{Lag}\\
\midrule
\texttt{tair}  & $+0.845$ & $0$ h & $+0.826$ & $0$ h\\
\texttt{twall} & $+0.791$ & $0$ h & $+0.477$ & $-3$ h\\
\texttt{sr}    & $+0.629$ & $+1$ h & $+0.584$ & $+1$ h\\
\texttt{rh}    & $-0.611$ & $0$ h & $-0.673$ & $0$ h\\
\bottomrule
\end{tabular}
\end{table}

Air temperature peaks at zero lag on both views. Solar radiation peaks one hour ahead
of the response on both. Wall temperature peaks at zero on levels but at $-3$ hours on
differences --- the inclination moves \emph{before} the wall probe does, consistent
with the probe sitting in a slower thermal mass than the one governing the tilt.

The sub-hourly check on the native 20-minute record agrees and refines the radiation
figure: the differenced peak for \texttt{sr} sits at $+20$ minutes, for \texttt{tair}
at zero and for \texttt{twall} at $-3$ hours. Nothing important hides below the
analysis resolution.

\subsection{Thermal inertia}

A shift is a crude model of a delayed response: it assumes the structure reproduces the
driver exactly, only later. A wall integrates instead. The standard lumped-capacitance
model of that behaviour is a single-pole low-pass filter with time constant $\tau$,
%
\begin{equation}
T_{\text{eff}}(t) \;=\; T_{\text{eff}}(t-\Delta t)
\;+\; \frac{\Delta t}{\tau + \Delta t}
\bigl( T_{\text{driver}}(t) - T_{\text{eff}}(t-\Delta t) \bigr) ,
\end{equation}
%
whose step response is the familiar $1 - e^{-t/\tau}$. Scanning $\tau$ and refitting at
each value asks the physical question directly, and $\tau = 0$ is the instantaneous
assumption --- so the correction's premise sits inside the scan as a special case
rather than being argued about separately.

\begin{figure}[H]
\includegraphics{S1_31_st02_inertia_scan.png}
\caption{Explained variance and fitted slope against thermal time constant, full
signal, Window A.}
\label{fig:tau}
\end{figure}

\begin{figure}[H]
\includegraphics{S1_35_st02_inertia_scan_band.png}
\caption{The same scan after removing a centred weekly mean from both series, which
restricts the comparison to the diurnal band.}
\label{fig:tauband}
\end{figure}

\begin{table}[H]
\centering
\caption{Best time constant per driver. The band-limited scan removes a centred weekly
mean and so describes the daily response alone; the full-signal scan is dominated by
Window A's seasonal warming.}
\small
\begin{tabular}{@{}l R{1.7cm} R{1.9cm} R{1.9cm} R{1.7cm} R{1.9cm} R{1.9cm}@{}}
\toprule
& \multicolumn{3}{c}{\textbf{Full signal}} & \multicolumn{3}{c}{\textbf{Diurnal band}}\\
\cmidrule(lr){2-4}\cmidrule(lr){5-7}
\textbf{Driver} & \textbf{Best $\tau$} & \textbf{$R^2$, $\tau{=}0$} & \textbf{$R^2$, best}
& \textbf{Best $\tau$} & \textbf{$R^2$, $\tau{=}0$} & \textbf{$R^2$, best}\\
\midrule
\texttt{tair}  & 0 h & $0.713$ & $0.713$ & 0 h & $\mathbf{0.932}$ & $\mathbf{0.932}$\\
\texttt{twall} & 0 h & $0.625$ & $0.625$ & 0 h & $0.588$ & $0.588$\\
\texttt{sr}    & \textbf{4 h} & $0.346$ & $\mathbf{0.460}$ & \textbf{2 h} & $0.575$ & $\mathbf{0.731}$\\
\texttt{rh}    & 0 h & $0.373$ & $0.373$ & 0 h & $0.484$ & $0.484$\\
\bottomrule
\end{tabular}
\end{table}

\subsection{What the four tests say}

\begin{description}
\item[Radiation acts with a lag, and the lag is physical.] Filtering the radiometer
through a two-hour thermal lag lifts the explained variance of the diurnal band from
$0.575$ to $0.731$; on the full signal a four-hour constant lifts it from $0.346$ to
$0.460$. Both are large gains bought with one parameter, and both point the same way.
The mechanism of Section~\ref{sec:hypothesis} is doing what it should: flux arrives, a
mass integrates it, the structure moves afterwards. The one-hour peak in the
cross-correlation is the same fact seen through a cruder instrument.
\item[Air temperature acts without a lag.] Every scan returns $\tau = 0$ for
\texttt{tair}, and at the diurnal scale instantaneous air temperature explains
$93.2\,\%$ of the variance of the inclination. Adding inertia makes the fit
monotonically worse.
\item[These two findings are consistent, not contradictory.] Air temperature is itself
a lagged, integrated response to the same radiation. Using it as the driver supplies the
delay implicitly, which is why it needs no further filtering --- and why it serves so
well as a proxy in the legacy era, where no radiometer exists. What the correction gets
right about timing, it gets right by accident of the proxy it happens to use.
\item[None of this rescues the correction.] Its defect is magnitude, not timing, and no
choice of lag or time constant changes the results in Sections~\ref{sec:size}
to~\ref{sec:experiment}.
\end{description}

\noindent\textbf{The limit of this window.} Window A spans half an annual cycle, so it
says nothing about the seasonal response. The companion legacy study, with a full
calendar year, finds that at seasonal scale the inclination \emph{leads} air
temperature by about two months --- a relationship no causal thermal lag can produce,
since a filter can only delay a driver and never advance it. The daily answer given
here does not generalise to the year.

%% ==========================================================================
\section{How large is the thermal response really?}
\label{sec:size}

\begin{table}[H]
\centering
\caption{Univariate regression of the inclination on each driver over Window A, with
heteroskedasticity- and autocorrelation-consistent (Newey--West) standard errors. The
comparison value is the $5.0$ \mdegC\ implied by the documented coefficient.}
\begin{tabular}{@{}l R{2.4cm} R{3.4cm} R{1.8cm} R{1.6cm} R{1.8cm}@{}}
\toprule
\textbf{Driver} & \textbf{Slope} & \textbf{95\,\% CI} & \textbf{SE} & \textbf{$R^2$} & \textbf{$n$}\\
\midrule
\multicolumn{6}{@{}l}{\itshape Levels}\\
\texttt{tair}  & $1.6279$ & $1.4854$ -- $1.7705$ & $0.0727$ & $0.713$ & 3\,419\\
\texttt{twall} & $1.2476$ & $1.1443$ -- $1.3508$ & $0.0527$ & $0.625$ & 3\,419\\
\texttt{sr}    & $0.0214$ & $0.0201$ -- $0.0227$ & $0.0007$ & $0.346$ & 3\,415\\
\texttt{rh}    & $-0.4341$ & $-0.5111$ -- $-0.3571$ & $0.0393$ & $0.373$ & 3\,419\\[2pt]
\multicolumn{6}{@{}l}{\itshape First differences}\\
\texttt{tair}  & $2.1730$ & $2.0927$ -- $2.2533$ & $0.0410$ & $0.682$ & 3\,414\\
\texttt{twall} & $0.8394$ & $0.7821$ -- $0.8967$ & $0.0292$ & $0.155$ & 3\,414\\
\texttt{sr}    & $0.0075$ & $0.0066$ -- $0.0084$ & $0.0005$ & $0.158$ & 3\,407\\
\texttt{rh}    & $-0.4544$ & $-0.4838$ -- $-0.4249$ & $0.0150$ & $0.454$ & 3\,414\\
\bottomrule
\end{tabular}
\end{table}

The documented $5.0$ \mdegC\ is not inside the confidence interval and is not near it.
The differenced estimate, $2.17$ \mdegC, describes the hourly response rather than the
seasonal one and is the larger of the two, but it is still less than half the
documented value. No timescale available in this window produces $5$.

\begin{figure}[H]
\includegraphics{S1_12_st02_thermal_scatter.png}
\caption{Inclination against air temperature (left) and wall temperature (right) over
Window A, with the fitted slope and the slope implied by the documented coefficient.}
\label{fig:scatter}
\end{figure}

The documented line is roughly three times steeper than the data. It crosses the cloud
near $16$\,\degC\ and departs from it in both directions, under-predicting the
inclination in the cold half of the record and over-predicting it in the warm half by
tens of millidegrees. The scatter also shows a second, lower branch between roughly
$15$ and $25$\,\degC, displaced some $30$ mdeg below the main relation: the same air
temperature corresponded to a different inclination at different times, which is direct
evidence that the level relation is not a single-valued function of temperature.

%% ==========================================================================
\section{Acceptance tests}

\begin{table}[H]
\centering
\caption{Acceptance tests on Window A. The target for every row is a \emph{lower} value
after correction. Eight of nine fail.}
\begin{tabular}{@{}l R{2.4cm} R{2.8cm} R{2.2cm} l@{}}
\toprule
\textbf{Test} & \textbf{Raw} & \textbf{Compensated} & \textbf{Change} & \textbf{Verdict}\\
\midrule
Variance [mdeg$^2$]                       & $211.90$ & $709.34$ & $+234.7\,\%$ & \worse\\
Standard deviation [mdeg]                 & $14.557$ & $26.633$ & $+83.0\,\%$  & \worse\\
Mean diurnal amplitude [mdeg]             & $23.520$ & $26.946$ & $+14.6\,\%$  & \worse\\
$|r|$ with \texttt{tair}, levels          & $0.8446$ & $0.9562$ & $+13.2\,\%$  & \worse\\
$|r|$ with \texttt{tair}, differenced     & $0.8260$ & $0.8856$ & $+7.2\,\%$   & \worse\\
$|r|$ with \texttt{twall}, levels         & $0.7908$ & $0.8781$ & $+11.0\,\%$  & \worse\\
$|r|$ with \texttt{twall}, differenced    & $0.3938$ & $0.4266$ & $+8.3\,\%$   & \worse\\
$|r|$ with \texttt{sr}, levels            & $0.5880$ & $0.4784$ & $-18.6\,\%$  & \better\\
$|r|$ with \texttt{sr}, differenced       & $0.3978$ & $0.6237$ & $+56.8\,\%$  & \worse\\
\bottomrule
\end{tabular}
\end{table}

The single improvement is not a success. It is what happens when a series is dragged
past zero: the absolute correlation with a partially independent driver passes through
a minimum on its way to becoming large and negative.

\begin{table}[H]
\centering
\caption{Signed correlations before and after correction. Every driver reverses sign,
which is the signature of overcorrection.}
\begin{tabular}{@{}l R{2.6cm} R{3.2cm} c@{}}
\toprule
\textbf{Driver} & \textbf{Raw} & \textbf{Compensated} & \textbf{Sign reversed}\\
\midrule
\texttt{tair}  & $+0.8446$ & $-0.9562$ & yes\\
\texttt{twall} & $+0.7908$ & $-0.8781$ & yes\\
\texttt{sr}    & $+0.5880$ & $-0.4784$ & yes\\
\bottomrule
\end{tabular}
\end{table}

\begin{figure}[H]
\includegraphics{S1_15_st02_compensation_effect.png}
\caption{Effect of the documented correction on Window A. Top left, the full window
before and after; top right, the first week in detail; bottom left, the thermal
dependence; bottom right, the mean diurnal cycle.}
\label{fig:effect}
\end{figure}

Each panel says something different.

\begin{description}
\item[Full window.] The raw series rises through spring, peaks in June and turns down
in July, a total excursion of roughly $60$ mdeg. The compensated series does not
flatten --- it acquires a large downward drift, reaching about $-100$ mdeg by mid-July.
A correction that worked would have removed the seasonal excursion; this one replaced
it with a larger one of opposite sign.
\item[First week.] At weekly scale the raw and compensated traces are near mirror
images about their own means.
\item[Thermal dependence.] The raw cloud slopes upward at about $1.6$ \mdegC. The
compensated cloud slopes downward at roughly $-3.4$ \mdegC\ --- which is $1.6 - 5.0$,
exactly what Equation~\eqref{eq:comp} implies when the true slope is $1.6$. The
correction does not land near zero; it lands on the other side and further away.
\item[Mean diurnal cycle.] The raw profile swings about $20$ mdeg with a midday
maximum. The compensated profile swings slightly \emph{more} and has a midday
\emph{minimum}. The daily thermal signature was not removed; it was inverted and
amplified.
\end{description}

%% ==========================================================================
\section{Where the documented coefficient sits}

\begin{table}[H]
\centering
\caption{Selected points from the coefficient sweep over Window A.}
\begin{tabular}{@{}l R{3.0cm} R{3.4cm} R{3.0cm}@{}}
\toprule
& \textbf{Coefficient [\mdegC]} & \textbf{Residual variance [mdeg$^2$]} & \textbf{$|r|$ with \texttt{tair}}\\
\midrule
                        & $-1.0$ & $454.66$  & $0.931$\\
No compensation         & $0.0$  & $211.90$  & $0.845$\\
                        & $1.0$  & $83.22$   & $0.520$\\
                        & $1.4$  & $63.69$   & $0.216$\\
\textbf{Variance-minimising} & $\mathbf{1.6}$ & $\mathbf{60.77}$ & $\mathbf{0.027}$\\
                        & $2.0$  & $68.62$   & $0.339$\\
                        & $3.0$  & $168.11$  & $0.799$\\
                        & $4.0$  & $381.68$  & $0.917$\\
\textbf{Documented}     & $\mathbf{5.0}$ & $\mathbf{709.34}$ & $\mathbf{0.956}$\\
                        & $6.0$  & $1151.08$ & $0.973$\\
\bottomrule
\end{tabular}
\end{table}

\begin{figure}[H]
\includegraphics{S1_17_st02_coefficient_sweep.png}
\caption{Residual variance and residual thermal correlation against the compensation
coefficient, with the documented value marked.}
\label{fig:sweep}
\end{figure}

Three readings decide criterion 2. First, a sharp minimum exists: at $1.6$ \mdegC\ the
residual variance falls to $60.8$ mdeg$^2$, less than a third of the uncorrected
$211.9$, and the residual thermal correlation falls to $0.027$. At that coefficient the
correction does exactly what it claims --- the mechanism is sound and only the constant
is wrong.

Second, the documented value sits far past the minimum on the ascending branch:
$709.3$ mdeg$^2$, which is $3.3\times$ the untouched signal and $11.7\times$ the
achievable minimum. That is the difference between a suboptimal correction and a
harmful one.

Third, the variance-minimising coefficient $1.60$ and the independently fitted
regression slope $1.63$ agree to within two per cent, from two different estimators.

%% ==========================================================================
\section{The prediction experiment}
\label{sec:experiment}

\subsection{One hour ahead}

\begin{table}[H]
\centering
\caption{Window A, one-hour-ahead ridge forecast, all errors on the raw signal scale.
Training and test sizes are 2\,351 and 1\,008 samples without regressors, 2\,348 and
1\,007 with.}
\small
\begin{tabular}{@{}l R{2.4cm} R{2.4cm} R{2.4cm}@{}}
\toprule
\textbf{Configuration} & \textbf{MAE [mdeg]} & \textbf{RMSE [mdeg]} & \textbf{vs best}\\
\midrule
D $\cdot$ raw, no regressors          & $\mathbf{1.9246}$ & $3.0875$ & \textbf{best}\\
C $\cdot$ compensated, no regressors  & $3.3075$ & $5.0805$ & $+71.9\,\%$\\
A $\cdot$ raw + instantaneous         & $6.5422$ & $8.0476$ & $+239.9\,\%$\\
F $\cdot$ raw + inertia-filtered      & $6.8054$ & $8.4610$ & $+253.6\,\%$\\
E $\cdot$ raw + lag-optimised         & $6.8648$ & $8.4968$ & $+256.7\,\%$\\
B $\cdot$ compensated + instantaneous & $7.0425$ & $8.4657$ & $+265.9\,\%$\\
\bottomrule
\end{tabular}
\end{table}

\begin{figure}[H]
\includegraphics{S1_19_st02_experiment_h1.png}
\caption{One-hour-ahead forecast, ridge model.}
\label{fig:h1}
\end{figure}

The controlled contrast is decided here: compensating the target costs $71.9\,\%$ in
mean absolute error with the model form held fixed. Every configuration carrying
exogenous regressors does worse still, E and F included --- at a one-step horizon the
previous value carries almost all the information, and Window A's rising temperature
forces any regressor-using linear model to extrapolate outside its training range. This
is the confound the C-versus-D design exists to avoid.

\subsection{Twenty-four hours ahead}

\begin{table}[H]
\centering
\caption{Window A, 24-hour-ahead ridge forecast.}
\small
\begin{tabular}{@{}l R{2.4cm} R{2.4cm} R{2.4cm}@{}}
\toprule
\textbf{Configuration} & \textbf{MAE [mdeg]} & \textbf{RMSE [mdeg]} & \textbf{vs best}\\
\midrule
D $\cdot$ raw, no regressors          & $\mathbf{3.6394}$ & $5.0113$ & \textbf{best}\\
C $\cdot$ compensated, no regressors  & $7.6349$ & $9.2839$ & $+109.8\,\%$\\
B $\cdot$ compensated + instantaneous & $13.1930$ & $16.1338$ & $+262.5\,\%$\\
F $\cdot$ raw + inertia-filtered      & $13.3373$ & $16.1089$ & $+266.5\,\%$\\
E $\cdot$ raw + lag-optimised         & $13.5486$ & $16.2908$ & $+272.3\,\%$\\
A $\cdot$ raw + instantaneous         & $14.1391$ & $16.8804$ & $+288.5\,\%$\\
\bottomrule
\end{tabular}
\end{table}

\begin{figure}[H]
\includegraphics{S1_21_st02_experiment_h24.png}
\caption{Twenty-four-hour-ahead forecast, ridge model.}
\label{fig:h24}
\end{figure}

The compensation penalty grows to $+109.8\,\%$. Among the regressor-using
configurations the ordering now favours the lagged and inertia-filtered forms over the
instantaneous one --- F and E beat A by $5.7\,\%$ and $4.2\,\%$ --- the first sign that
the measured lag structure carries real information, though all three remain far behind
plain autoregression in this seasonally extrapolating window.

\subsection{Out-of-period check on Window B}

\begin{table}[H]
\centering
\caption{Window B, 24-hour-ahead ridge forecast. Different season, different year, far
side of a 103-day outage; 525 training and 225 test samples.}
\small
\begin{tabular}{@{}l R{2.4cm} R{2.4cm} R{2.4cm}@{}}
\toprule
\textbf{Configuration} & \textbf{MAE [mdeg]} & \textbf{RMSE [mdeg]} & \textbf{vs best}\\
\midrule
D $\cdot$ raw, no regressors          & $\mathbf{3.3318}$ & $4.2758$ & \textbf{best}\\
E $\cdot$ raw + lag-optimised         & $5.0081$ & $5.9162$ & $+50.3\,\%$\\
F $\cdot$ raw + inertia-filtered      & $5.2299$ & $6.0995$ & $+57.0\,\%$\\
A $\cdot$ raw + instantaneous         & $5.7027$ & $6.5291$ & $+71.2\,\%$\\
B $\cdot$ compensated + instantaneous & $5.9197$ & $7.2911$ & $+77.7\,\%$\\
C $\cdot$ compensated, no regressors  & $7.3078$ & $9.3234$ & $+119.3\,\%$\\
\bottomrule
\end{tabular}
\end{table}

In Window B the seasonal-extrapolation confound is absent --- the window is midsummer
throughout --- and two things follow. The compensated configurations are worst, and
among the regressor-using ones both lagged forms beat the instantaneous one, E by
$12\,\%$ and F by $8\,\%$ relative to A. Where the model is not being asked to
extrapolate, respecting the lag structure pays.

\subsection{The controlled contrast, collected}

\begin{table}[H]
\centering
\caption{C against D across both windows and both horizons. Positive penalty means the
compensated target was harder to forecast.}
\begin{tabular}{@{}c l R{2.6cm} R{3.0cm} R{2.6cm}@{}}
\toprule
\textbf{Window} & \textbf{Horizon} & \textbf{MAE raw (D)} & \textbf{MAE compensated (C)} & \textbf{Penalty}\\
\midrule
A & 1 h  & $1.9246$ & $3.3075$ & $+71.9\,\%$\\
A & 24 h & $3.6394$ & $7.6349$ & $+109.8\,\%$\\
B & 1 h  & $3.5475$ & $5.9657$ & $+68.2\,\%$\\
B & 24 h & $3.3318$ & $7.3078$ & $+119.3\,\%$\\
\bottomrule
\end{tabular}
\end{table}

Four measurements, two windows, two horizons, eleven months apart. The compensated
target is harder to forecast in every one of them.

\subsection{Confirmation under NeuralProphet}
\label{sec:np}

The ridge model establishes the ordering cheaply and transparently. This repeats it
under the model the project deploys --- 24 autoregressive lags, 20 epochs,
chronological 70/30 split, one hour ahead, fixed random seed.

\begin{table}[H]
\centering
\caption{Window A, one hour ahead under NeuralProphet, 1\,030 scored points per
configuration, all errors on the raw signal scale.}
\small
\begin{tabular}{@{}l R{2.4cm} R{2.4cm} R{2.4cm}@{}}
\toprule
\textbf{Configuration} & \textbf{MAE [mdeg]} & \textbf{RMSE [mdeg]} & \textbf{vs best}\\
\midrule
A $\cdot$ raw + instantaneous         & $\mathbf{2.2249}$ & $2.8271$ & \textbf{best}\\
E $\cdot$ raw + lag-optimised         & $2.2259$ & $2.8299$ & $+0.0\,\%$\\
D $\cdot$ raw, no regressors          & $2.2381$ & $2.8688$ & $+0.6\,\%$\\
F $\cdot$ raw + inertia-filtered      & $2.3124$ & $2.9605$ & $+3.9\,\%$\\
C $\cdot$ compensated, no regressors  & $3.8803$ & $5.0989$ & $+74.4\,\%$\\
B $\cdot$ compensated + instantaneous & $4.3028$ & $5.6304$ & $+93.4\,\%$\\
\bottomrule
\end{tabular}
\end{table}

\begin{figure}[H]
\includegraphics{S1_24_st02_experiment_neuralprophet.png}
\caption{One-hour-ahead forecast under NeuralProphet.}
\label{fig:np}
\end{figure}

The result separates cleanly into two groups. \textbf{The four configurations built on
the raw signal are indistinguishable}: they span $2.22$ to $2.31$ mdeg, under four per
cent, which is within the run-to-run variability of training this model and should be
read as a tie. \textbf{The two built on the compensated signal are far worse}, by
$74\,\%$ and $93\,\%$. That gap is an order of magnitude larger than the spread among
the raw configurations, and it is the finding that survives.

Training is stochastic, so the notebook fixes the seed. Seeding makes the ordering
reproducible; it does not make small differences meaningful, and no claim in this report
rests on one.

%% ==========================================================================
\section{Verdict}

\panel{Three criteria, three failures}{%
\textbf{1. Thermal dependence is not reduced.} The absolute correlation with air
temperature rises from $0.845$ to $0.956$ and reverses sign.\par\medskip
\textbf{2. The coefficient is not near the variance-minimising value.} $5.00$ against
$1.60$ \mdegC, inflating residual variance to $3.3\times$ the uncorrected
signal.\par\medskip
\textbf{3. The forecast does not improve.} The controlled-contrast penalty ranges from
$+68.2\,\%$ to $+119.3\,\%$ across two windows and two horizons, and is $+74.4\,\%$
under NeuralProphet. No window, horizon or model favoured the compensated target.}

\subsection{Recommendation}

\textbf{Remove the fixed compensation from Notebook 00 and let the model estimate the
thermal response.} Every raw configuration outperforms every compensated one under both
models tested. The environmental response is exactly what the grey-box decomposition
exists to identify; subtracting a guess at it beforehand discards information the
decomposition needs and injects an artefact it must then remove.

If a compensated series is wanted for presentation, the coefficient must be
\textbf{fitted per station}, not fixed. At $1.60$ \mdegC\ the correction does what it
claims. But the legacy stations gave $7.21$, $1.33$ and $9.66$ \mdegC\ --- a sevenfold
spread across three nominally identical instruments on one wall --- so no single
constant can serve all of them, and the constant currently in the code serves none.

\subsection{On modelling the thermal response instead}

The lag analysis says what should replace the correction. Air temperature can be used
instantaneously, being a proxy that already carries the delay. Solar radiation should be
used through a thermal lag of two to four hours rather than raw: that single parameter
is worth $0.16$ in explained variance at the diurnal scale. Wall temperature, as
recorded by this probe, lags the tilt rather than leading it, and is the least useful of
the three as a predictor whatever its value as a diagnostic.

%% ==========================================================================
\section{Caveats}

\begin{itemize}
\item Window A spans one half of an annual cycle. Nothing here constrains the seasonal
response, and the companion legacy study shows the seasonal picture is qualitatively
different from the daily one.
\item The band-limited inertia scan removes a centred weekly mean, so it describes
variation faster than a week. The window width is a modelling choice; a different width
moves the numbers, though the ordering of the drivers is stable.
\item The levels slope includes the seasonal trend and the differenced slope describes
the hourly timescale. Both are reported, they differ by a factor of $1.3$, and neither
is near $5.0$.
\item The regressor-using ridge configurations are handicapped in Window A by seasonal
extrapolation, which is why the verdict rests on C against D.
\item Differences below about five per cent between NeuralProphet configurations are
within training variability and are treated as ties.
\item The manufacturer evidence in Section~\ref{sec:manufacturer} comes from a separate
review of the supplied documentation, not from this notebook. Every number in
Sections~2 to~8 stands without it.
\end{itemize}

%% ==========================================================================
\appendix
\section{Reproduction}

\begin{verbatim}
conda activate neuralprophet_env
cd studies/thermal_compensation
jupytext --to ipynb thermal_compensation_study.py
jupyter nbconvert --to notebook --execute --inplace \
        thermal_compensation_study.ipynb
\end{verbatim}

\noindent Edit the paired \texttt{.py} file, never the \texttt{.ipynb}. The first run
copies 416 files out of the read-only Google Drive archive, which takes roughly forty
minutes at the observed streaming rate; later runs reuse \texttt{.cache/}. The archive
is never written to.

Setting \texttt{CONTEXT = 'paper'} in the parameter cell re-exports every figure at
manuscript column size: seaborn rescales the typography and the shared canvas scale
rescales the figures, so the exported figure is the same figure rather than a different
one.

\section{Output inventory}

Every figure and table in this report is regenerated by executing the notebook. The
\texttt{outputs/} directory is gitignored.

\begin{table}[H]
\centering
\small
\caption{Artefacts produced by the notebook. Identifiers are stable names, not an
ordering.}
\begin{tabular}{@{}l L{9.2cm} c@{}}
\toprule
\textbf{Identifier} & \textbf{Contents} & \textbf{In this report}\\
\midrule
S1\_01 & Channel coverage & Table 1\\
S1\_02 & Window A overview & Fig.~\ref{fig:ovA}\\
S1\_03 & Window B overview & Fig.~\ref{fig:ovB}\\
S1\_04 & Mean diurnal profiles & Fig.~\ref{fig:diurnal}\\
S1\_05, S1\_06 & Correlation matrices, levels and differences & Table 4\\
S1\_07 & Correlation heatmaps & Fig.~\ref{fig:heat}\\
S1\_08 & Cross-correlation peaks, differences & Table 5\\
S1\_09 & Lag structure, levels and differences & Fig.~\ref{fig:ccf}\\
S1\_10, S1\_11 & Regression slopes, levels and differences & Table 7\\
S1\_12 & Fitted versus documented thermal slope & Fig.~\ref{fig:scatter}\\
S1\_13 & Acceptance tests & Table 8\\
S1\_14 & Signed correlation check & Table 9\\
S1\_15 & Effect of the documented compensation & Fig.~\ref{fig:effect}\\
S1\_16 & Coefficient sweep & Table 10\\
S1\_17 & Coefficient sweep figure & Fig.~\ref{fig:sweep}\\
S1\_18, S1\_20 & Ridge results, 1 h and 24 h & Tables 11, 12\\
S1\_19, S1\_21 & Experiment figures, 1 h and 24 h & Figs.~\ref{fig:h1}, \ref{fig:h24}\\
S1\_22 & Ridge results, Window B & Table 13\\
S1\_23 & NeuralProphet results & Table 15\\
S1\_24 & NeuralProphet experiment figure & Fig.~\ref{fig:np}\\
S1\_25 & Verdict & Section 8\\
S1\_26 & Controlled contrast, collected & Table 14\\
S1\_27 & Cross-correlation peaks, levels & Table 5\\
S1\_28 & Cross-correlation peaks, native 20-minute record & Section~\ref{sec:lag}\\
S1\_29, S1\_30 & Inertia scan and best time constants, full signal & Table 6\\
S1\_31 & Inertia scan figure, full signal & Fig.~\ref{fig:tau}\\
S1\_32 & Lag summary & Section~\ref{sec:lag}\\
S1\_33, S1\_34 & Inertia scan and best time constants, diurnal band & Table 6\\
S1\_35 & Inertia scan figure, diurnal band & Fig.~\ref{fig:tauband}\\
\bottomrule
\end{tabular}
\end{table}

\end{document}
